\scanpage{7}
Considérons le cas d'un \uline{automorphisme} $f$, et voyons quel est le passage au \unsure{contragrédient|contragredient}
\[
\begin{aligned}
P_{\check f}(t)&=\det(1-\check f t)=\prod(1-\lambda_i^{-1}\deleted{\unsure{f}}t)\\
&\deleted{=\prod\lambda_i^{-1}(\unsure{\lambda_i-t|\bar\lambda_i-t})}\\
&=\prod_i(-\lambda_i^{-1}t)(1-\lambda_i t^{-1})\\
&=\frac{(-1)^nt^n}{\prod\lambda_i}\deleted{\unsure{\det}}\prod(1-\lambda_i t^{-1})\\
&=\frac{(-1)^nt^n}{\det f}P_f(\unsure{t^{-1}|-t^{-1}})\qquad\text{donc}
\end{aligned}
\]
\[
\deleted{P_{\check f}(t)=\deleted{\frac{(-1)^nt^n}{\det f}}\unsure{\text{``}\varepsilon(f)\text{''}|{\varepsilon(f)}}t^n P_f(t^{-1})\qquad\varepsilon(f)=}
\]
\[
\boxed{P_{\check f}(t)=(-t)^{\varepsilon(f)}\det f^{-1}\ P_f(t^{-1})}\qquad\text{i.e}
\]
\[
\boxed{L_{\check f}(t)=(-t)^{-\varepsilon(f)}\det f\ L_f(t^{-1})}
\]
Bien entendu, le deuxième \unsure{mb} n'a pas de sens quand on remplace $P_f^{\unsure{=|-}}(t)$ par une série formelle quelconque : il est essentiel de se borner aux séries formelles qui sont \unsure{produits|justement} alternés finis de \uline{polynômes} [à terme constant 1, à terme dominant inversible -- cette dernière condition n'étant utile que si $k$ est un anneau général \ldots].
% archive sheet number: 6

